\begin{statementbox}
	\textbf{\textcolor{CStatement}{条件}}
	\begin{description}[style=nextline, leftmargin=2.8em, labelsep=0.5em, itemsep=0.35em]
		\item[\textbf{\textcolor{CStatement}{条件1（几何场景）：}}]
			给定直线$l$，两定点$A,B$在$l$的同侧，动点$P$在$l$上运动。
		\item[\textbf{\textcolor{CStatement}{条件2（参数构造）：}}]
			常数$k\in(0,1)$；过点$A$作射线$AN$，使$AN$与$l$的夹角$\alpha$满足$\sin\alpha=k$，且射线$AN$与点$B$分居直线$l$的两侧。
	\end{description}
	\textbf{\textcolor{CStatement}{结论}}
	\begin{description}[style=nextline, leftmargin=2.8em, labelsep=0.5em, itemsep=0.45em]
		\item[\textbf{\textcolor{CStatement}{结论一（构造转化）：}}]
			\textbf{\textcolor{CStatement}{直观描述：}}
			构造满足$\sin\alpha=k$的射线$AN$，将$k\cdot PA$转化为$P$到$AN$的垂线段长，则$k\cdot PA+PB$的最小值即为点$B$到$AN$的距离$BH$，当$P$为$BH$与$l$的交点时取等。
			\[
				\min_{P\in l}(k\cdot PA+PB)=BH.
			\]
		\item[\textbf{\textcolor{CStatement}{推广结论（系数调整）：}}]
			\textbf{\textcolor{CStatement}{直观描述：}}
			当$k>1$时，提取公因数使括号内系数小于1，再应用原结论。
			\[
				k\cdot PA+PB = k\left(PA+\frac{1}{k}PB\right),\quad \frac{1}{k}\in(0,1).
			\]
	\end{description}

\end{statementbox}
